% !TeX spellcheck = es_ES
\documentclass[letterpaper,11pt]{article}
\usepackage[utf8]{inputenc}
\usepackage[T1]{fontenc}
\usepackage[spanish]{babel}
\decimalpoint % separador decimal con punto, como en las tablas de la fuente original
\usepackage{lmodern,microtype}
\usepackage{amsmath,amssymb,amsfonts}
\usepackage{geometry}
\geometry{margin=2.54cm}
\usepackage{setspace}
\setstretch{1.15}
\usepackage{enumitem}
\usepackage{booktabs}
\usepackage{color}
\usepackage{hyperref}
\definecolor{ntr}{rgb}{0.8,0,0}
\hypersetup{colorlinks=true,linkcolor=ntr,citecolor=ntr,urlcolor=black}

% Portada (Luis Cabral): titulo \LARGE sans-serif alineado a la izquierda, sin numero de pagina
\makeatletter
    \renewcommand{\@maketitle}{%
    \newpage\parindent0in\null\vskip2em%
    {\LARGE\textbf{\textsf{\@title}}\par\vskip2em}%
    \@author\par\vskip2em%
    \@date%
    \par
    \thispagestyle{empty}\setcounter{page}{1}%
    }%
\makeatother

\DeclareMathOperator{\Var}{Var}
\DeclareMathOperator{\Cov}{Cov}
\DeclareMathOperator{\Corr}{Corr}
\DeclareMathOperator{\E}{\mathbb{E}}
\DeclareMathOperator{\N}{Normal}
\DeclareMathOperator{\LP}{\mathbb{L}}
\newcommand\independent{\protect\mathpalette{\protect\independenT}{\perp}}
\def\independenT#1#2{\mathrel{\rlap{$#1#2$}\mkern2mu{#1#2}}}
% Teoria asintotica
\newcommand{\SR}{\mathbb{R}}
\DeclareMathOperator{\plim}{plim}
\DeclareMathOperator{\Avar}{AVar}
\newcommand{\pto}{\stackrel{p}{\longrightarrow}}
\newcommand{\dto}{\stackrel{d}{\longrightarrow}}
\newcommand{\adistr}{\stackrel{a}{\sim}}
\newcommand{\epsi}{\varepsilon}


\title{Análisis Econométrico \\ Ayudantía 4}
\author{Magíster en Economía, Universidad Alberto Hurtado}
\date{}

\begin{document}
\maketitle

\noindent Esta ayudantía comienza con un control breve (\textbf{Control 2}, 15 minutos, sin
apuntes) sobre el contenido de la Tarea 2. Luego revisamos en conjunto los ejercicios 1, 2 y 4
de la Tarea 2 --- se reproducen aquí para tenerlos a mano durante la revisión.

\vspace{1em}

\noindent\textbf{Ejercicio 1 de la Tarea 2 (20 puntos). Bernoulli.}

Dada una muestra i.i.d.\ de tamaño $N$ de una variable aleatoria $y$ con distribución Bernoulli
con probabilidad de éxito $p$. Es decir,
\begin{align*}
y = \begin{cases}
1 & \text{con probabilidad } p,\\
0 & \text{con probabilidad } 1-p.
\end{cases}
\end{align*}
En lo que sigue puede utilizar que $\E(y)=p$ y $\Var(y)=p(1-p)$.
\begin{enumerate}
    \item Demuestre que la media muestral $\bar y=\frac{1}{N}\sum_i y_i$ es un estimador
    consistente de $p$.
    \item Demuestre que $\bar y(1-\bar y)$ es un estimador consistente de $\Var(y)=p(1-p)$.
    \item Obtenga la distribución asintótica de $\sqrt{N}(\bar y-p)$. Obtenga una distribución
    aproximada de $\bar y$ para muestras grandes.
    \item Obtenga la distribución asintótica de
    $\sqrt{N}\frac{(\bar y-p)}{\sqrt{\bar y(1-\bar y)}}$. Explique por qué este resultado es más
    útil en la práctica que el de la parte (c).
\end{enumerate}

\vspace{1em}

\noindent\textbf{Ejercicio 2 de la Tarea 2 (20 puntos). Covarianza muestral.}

La covarianza muestral se define como
$S_{xy}=\frac{1}{N}\sum_{i=1}^N (x_i-\bar x)(y_i-\bar y)$ y la covarianza poblacional como
$\sigma_{xy}=\E(xy)-\E(x)\E(y)$. Sean $\mu_x=\E(x)$ y $\mu_y=\E(y)$.
\begin{enumerate}
    \item Demuestre que la covarianza muestral se puede escribir como
    \begin{align*}
    S_{xy}=\frac{1}{N}\sum_{i=1}^N (x_i-\mu_x)(y_i-\mu_y)-(\bar x-\mu_x)(\bar y-\mu_y).
    \end{align*}
    \item Demuestre que la covarianza muestral es un estimador consistente de la covarianza
    poblacional, $S_{xy}\pto \sigma_{xy}$.
    \item Demuestre que
    \begin{align*}
    \sqrt{N}(S_{xy}-\sigma_{xy})\dto \N\left(0,\ \E[(x-\mu_x)^2(y-\mu_y)^2]-\sigma_{xy}^2\right).
    \end{align*}
    Pista: parta de la descomposición de la parte (a), aplique el teorema central del límite al
    primer término y el teorema de Slutsky para mostrar que el segundo término es
    asintóticamente despreciable.
\end{enumerate}

\vspace{1em}

\noindent\textbf{Ejercicio 4 de la Tarea 2 (20 puntos). Método delta: cociente.}

$\hat\theta_N=(\hat\theta_{1N},\hat\theta_{2N})'$ es un estimador consistente y
$\sqrt{N}$-asintóticamente normal del vector $\theta=(\theta_1,\theta_2)'$ donde
$\theta_2\neq 0$. Defina $\hat\gamma_N=\hat\theta_{1N}/\hat\theta_{2N}$ como un estimador de
$\gamma=\theta_1/\theta_2$.
\begin{enumerate}
    \item Demuestre que $\hat\gamma_N$ es un estimador consistente de $\gamma$.
    \item Encuentre la varianza asintótica de $\hat\gamma_N$ ($\Avar(\hat\gamma_N)$) en función
    de $\theta$ y la varianza asintótica de $\hat\theta_N$ ($\Avar(\hat\theta_N)$).\\
    Pista: primero encuentre la distribución asintótica de $\sqrt{N}(\hat\gamma_N-\gamma)$.
    \item Suponga que se obtiene de una muestra aleatoria que $\hat\theta_N=(-4,2)'$ y una
    estimación de $\Avar(\hat\theta_N)$ es
    \begin{align*}
    \begin{pmatrix} 1 & -0{,}1 \\ -0{,}1 & 0{,}25 \end{pmatrix}.
    \end{align*}
    Calcule $\hat\gamma_N$ y su error estándar asintótico $se(\hat\gamma_N)$.
\end{enumerate}

\vspace{2em}
\noindent\rule{\textwidth}{0.4pt}

\section*{Material adicional (opcional)}

\noindent Práctica extra, no obligatoria y no cubierta en esta ayudantía. Es muy parecido al
ejercicio 4.

\vspace{1em}

\noindent\textbf{Ejercicio 3 de la Tarea 2 (20 puntos). Método delta: $\log(\hat\theta_N)$.}

$\hat\theta_N$ es un estimador consistente y $\sqrt{N}$-asintóticamente normal del parámetro
$\theta>0$. Defina $\hat\gamma_N=\log(\hat\theta_N)$ como un estimador de $\gamma=\log(\theta)$.
\begin{enumerate}
    \item Demuestre que $\hat\gamma_N$ es un estimador consistente de $\gamma$.
    \item Suponga que la varianza asintótica de $\sqrt{N}(\hat\theta_N-\theta)$ es $V$.
    Encuentre la distribución asintótica de $\sqrt{N}(\hat\gamma_N-\gamma)$.
    \item Encuentre la distribución asintótica de $\hat\gamma_N$. Suponga que se obtiene de una
    muestra aleatoria que $\hat\theta_N=4$ y $se(\hat\theta_N)=2$. Calcule $\hat\gamma_N$ y su
    error estándar asintótico $se(\hat\gamma_N)$.
\end{enumerate}

\vspace{1em}

\end{document}
